\begin{lemma}[最优轨道在容量弧上的延续]
\label{lem:no-early-boundary-exit}
在定理\ref{thm:verification}的条件下，设 $q\in\Uad(x_0)$ 为最优控制，
其状态轨道为 $x_q=(s_q,i_q)$，首次进入安全集的时间为 $\tau_q$。
若某个 $t_0\in[0,\tau_q)$ 满足
\[
x_q(t_0)=(s_*,K),\qquad s_*>s_B,
\]
令
\[
t_B:=t_0+T_B(s_*,s_B)
=t_0+\frac1{cK}\ln\frac{s_*-r}{s_B-r}.
\]
则 $t_B<\tau_q$，并且
\begin{equation}
i_q(t)=K,\qquad
s_q(t)=r+(s_*-r)\e^{-cK(t-t_0)}
\quad(t_0\le t\le t_B).
\label{eq:optimal-boundary-continuation}
\end{equation}
在 $(t_0,t_B)$ 上，$q(t)=q_B(s_q(t))$ 几乎处处。
特别地，最优轨道不能在 $s>s_B$ 时离开容量边界。
\end{lemma}

\begin{proof}
先证明局部结论。
由 $\Phi(s_*,K)>\Phi_B$ 及连续性，
可取 $(s_*,K)$ 在 $\Dcal$ 中的一个相对邻域 $O$，
使其中 $s>s_B$、$\Phi>\Phi_B$，且
\[
O\cap\{i<K\}\subset\mathcal W_K.
\]
轨道连续性给出 $\delta>0$，使
$t_0+\delta<\min\{\tau_q,t_B\}$ 且
$x_q([t_0,t_0+\delta])\subset O$。

若开时间集合
\[
E:=\{t\in(t_0,t_0+\delta):i_q(t)<K\}
\]
非空，取其任一连通分支 $(a,b)$。
由连续性及 $i_q(t_0)=K$，有 $a\ge t_0$ 和 $i_q(a)=K$。
在 $(a,b)$ 上轨道属于 $\mathcal W_K$。
推论\ref{cor:verification-equality}与
引理\ref{lem:waiting-HJB}因此给出 $q=0$ 几乎处处于 $(a,b)$。
于是对每个 $t\in(a,b)$，
\[
i_q(t)-K
=\int_a^t pc\bigl(s_q(v)-h\bigr)i_q(v)\dd v>0,
\]
因为 $s_q(v)>s_B>h$ 且 $i_q(v)>0$。
这与 $i_q<K$ 矛盾，故 $E$ 为空。
因此 $i_q=K$ 在 $[t_0,t_0+\delta]$ 上成立。

在任何保持 $i_q=K$ 的时间区间上，
引理\ref{lem:ac-level-set}与状态方程给出
\[
q(t)=1-\frac{h}{s_q(t)},\qquad
\dot s_q(t)=-cK(s_q(t)-r)
\]
几乎处处成立，从而得到 \eqref{eq:optimal-boundary-continuation} 中的显式解。

为将局部结论延续至 $t_B$，令
\[
t_*:=\sup\{T\in[t_0,t_B]:
 i_q(t)=K\text{ 对所有 }t\in[t_0,T]\text{ 成立}\}.
\]
由局部结论，$t_*>t_0$；由连续性，$i_q=K$ 在 $[t_0,t_*]$ 上成立。
若 $t_*<t_B$，上述显式解给出 $s_q(t_*)>s_B$。
同时，该段轨道始终位于安全集外，故 $t_*<\tau_q$。
在 $t_*$ 处重新应用局部结论，可将 $i_q=K$ 的区间继续向右延长，
与 $t_*$ 的定义矛盾。
因此 $t_*=t_B$，所述状态公式与控制公式在整个容量阶段成立。
最后，$x_q(t_B)=(s_B,K)$，且
$\Phi(s_B,K)=\Phi_B>\Phi_A$，所以 $t_B<\tau_q$。
\end{proof}
